An old planetary nebula, with a white dwarf (WD) in its center, is located
50\,pc away from Earth. Exactly in the same direction, but behind the
nebula, lies another WD, identical to the first, but located at 150\,pc from
the Earth. Consider that the two WDs have absolute bolometric magnitude
$+14.2$ and intrinsic color indexes $B-V = 0.300$ and $U-V = 0.330$.
Extinction occurs in the interstellar medium and in the planetary nebula.

When we measure the color indices for the closer WD (the one who lies at the
center of the nebula), we find the values $B-V = 0.327$ and $U-B = 0.038$.
In this part of the Galaxy, the interstellar extinction rates are 1.50, 1.23
and 1.00 magnitudes per kiloparsec for the filters $U$, $B$ and $V$,
respectively. Calculate the color indices as they would be measured for the
second star.
